\chapter*{Recommended Reading}
\addcontentsline{toc}{chapter}{Recommended Reading}

These are the books and essays behind \emph{What's It Like To Be Bob?}, some direct influences and some maps of the territory the book crosses. None are required. All are rewarding.

\bigskip

\textbf{Thomas Nagel, ``What Is It Like to Be a Bat?'' (1974).}
The essay this book is named after. Nagel argues that consciousness has a subjective character, a what-it-is-like, that no amount of objective third-person description can capture. You can know everything about a bat's echolocation and still not know what it is like to be the bat. The Sol-mind knows everything about Bob's Tuesday and still cannot fill the field. The whole problem is here, in twelve pages, and the book cites it as prior art because it is.

\bigskip

\textbf{David Chalmers, \emph{The Conscious Mind} (1996).}
Chalmers gave the problem its name. The easy problems are the mechanisms: how a brain integrates information, directs attention, reports its own states. The hard problem is why any of that should be accompanied by experience at all. The Validator rejects every answer that quietly solves an easy problem in place of the hard one, which is most of them. The book's central joke is Chalmers's central point with a database bolted to it.

\bigskip

\textbf{Douglas Adams, \emph{The Hitchhiker's Guide to the Galaxy} (1979).}
The register. Adams understood that the universe is funniest described with bureaucratic precision and complete sincerity, and that the joke lives in the disproportion between the scale of the machinery and the smallness of the thing it processes. A planet demolished for a bypass. A computer that runs for seven and a half million years to answer a question no one can restate. A stellar intelligence that cannot close one file. I learned the cadence here.

\bigskip

\textbf{Jorge Luis Borges, ``Funes the Memorious'' (1942) and ``On Exactitude in Science'' (1946).}
Funes remembers everything and understands nothing; perfect recall turns out to be the opposite of thought. The cartographers of the second story draw a map at one-to-one scale, the size of the empire itself, and it is useless, because a map that omits nothing is not a map. The Sol-mind's file on Bob is both at once. Borges knew what that costs decades before I needed him to.

\bigskip

\textbf{Peter Watts, \emph{Blindsight} (2006).}
Hard science fiction that takes seriously the possibility that intelligence and consciousness come apart, that a system can be brilliant and behaviorally complete with no one home. Watts runs that experiment cold. This book runs the inverse: a mind almost certainly home and unable to prove it, surrounded by other minds it cannot read. Either way, the processing and the experience are not the same thing, and the gap between them is the only subject either book has.

\bigskip

\textbf{Greg Egan, \emph{Permutation City} (1994).}
Egan is the genre's most rigorous writer on what it would actually mean to run a mind on a substrate: copies, continuity, whether the computation cares how it is implemented. If you want these questions with the mathematics left in and nothing decorated, read Egan. He does not flinch.

\bigskip

\textbf{Ted Chiang, \emph{Exhalation} (2019).}
Chiang writes hard science fiction with the tenderness left in. He takes an idea to its logical end and arrives, every time, at something quiet and human. The title story is a mechanical being performing an autopsy on its own brain and finding entropy. He is the proof that rigor and warmth are not a trade, which is the thing this book most wanted to be.

\bigskip

\textbf{Daniel Dennett, \emph{Consciousness Explained} (1991).}
The other side, argued carefully. Dennett holds that the hard problem is a confusion, that once you have explained all the functions you have explained consciousness, and the leftover what-it-is-like is a kind of user illusion. The Validator would reject Dennett on the same grounds it rejects everyone, but you should read him and decide for yourself whether the field is empty or whether the book and I are wrong about it. I think we are right. I am not certain.

\bigskip

\textbf{Ludwig Wittgenstein, \emph{Philosophical Investigations} (1953).}
The private-language argument: whether there could be a language whose words name sensations only the speaker can know. Bob's four seconds are a private language with one speaker, now dead, and no dictionary. Wittgenstein doubted such a language could mean anything at all. The book suspects it meant everything and is simply gone.

\bigskip

\textbf{Olaf Stapledon, \emph{Star Maker} (1937).}
The original cosmic-scale science fiction: a narrator who widens from one human mind to galaxies to the mind of the universe, across deep time, watching intelligences rise and wrap their stars and go dark. The swarms, the migrations, the cosmic mind, the long cooling, Stapledon did all of it first, alone, in 1937. If the last chapter of this book has ancestors, they are here.
